\documentclass[10pt,a4paper]{amsart}
\usepackage{amsmath,amssymb,amsthm,graphicx}
\usepackage{verbatim}
\newenvironment{reference}{\comment}{\endcomment}
%\newenvironment{reference}{}{}
\newenvironment{slogan}{\comment}{\endcomment}
\newenvironment{history}{\comment}{\endcomment}
% For commutative diagrams we use Xy-pic
\usepackage[all]{xy}

% We use 2cell for 2-commutative diagrams.
\xyoption{2cell}
% Theorem environments.
%
\theoremstyle{plain}
\newtheorem{theorem}[subsection]{Theorem}
\newtheorem{proposition}[subsection]{Proposition}
\newtheorem{lemma}[subsection]{Lemma}

\theoremstyle{definition}
\newtheorem{definition}[subsection]{Definition}
\newtheorem{example}[subsection]{Example}
\newtheorem{exercise}[subsection]{Exercise}
\newtheorem{situation}[subsection]{Situation}

\theoremstyle{remark}
\newtheorem{remark}[subsection]{Remark}
\newtheorem{remarks}[subsection]{Remarks}

\numberwithin{equation}{subsection}

% Macros
%
\def\lim{\mathop{\mathrm{lim}}\nolimits}
\def\colim{\mathop{\mathrm{colim}}\nolimits}
\def\Spec{\mathop{\mathrm{Spec}}}
\def\Hom{\mathop{\mathrm{Hom}}\nolimits}
\def\Ext{\mathop{\mathrm{Ext}}\nolimits}
\def\SheafHom{\mathop{\mathcal{H}\!\mathit{om}}\nolimits}
\def\SheafExt{\mathop{\mathcal{E}\!\mathit{xt}}\nolimits}
\def\Sch{\mathit{Sch}}
\def\Mor{\mathop{\mathrm{Mor}}\nolimits}
\def\Ob{\mathop{\mathrm{Ob}}\nolimits}
\def\Sh{\mathop{\mathit{Sh}}\nolimits}
\def\NL{\mathop{N\!L}\nolimits}
\def\CH{\mathop{\mathrm{CH}}\nolimits}
\def\proetale{{pro\text{-}\acute{e}tale}}
\def\etale{{\acute{e}tale}}
\def\QCoh{\mathit{QCoh}}
\def\Ker{\mathop{\mathrm{Ker}}}
\def\Im{\mathop{\mathrm{Im}}}
\def\Coker{\mathop{\mathrm{Coker}}}
\def\Coim{\mathop{\mathrm{Coim}}}

% Boxtimes
%
\DeclareMathSymbol{\boxtimes}{\mathbin}{AMSa}{"02}

%
% Macros for moduli stacks/spaces
%
\def\QCohstack{\mathcal{QC}\!\mathit{oh}}
\def\Cohstack{\mathcal{C}\!\mathit{oh}}
\def\Spacesstack{\mathcal{S}\!\mathit{paces}}
\def\Quotfunctor{\mathrm{Quot}}
\def\Hilbfunctor{\mathrm{Hilb}}
\def\Curvesstack{\mathcal{C}\!\mathit{urves}}
\def\Polarizedstack{\mathcal{P}\!\mathit{olarized}}
\def\Complexesstack{\mathcal{C}\!\mathit{omplexes}}
% \Pic is the operator that assigns to X its picard group, usage \Pic(X)
% \Picardstack_{X/B} denotes the Picard stack of X over B
% \Picardfunctor_{X/B} denotes the Picard functor of X over B
\def\Pic{\mathop{\mathrm{Pic}}\nolimits}
\def\Picardstack{\mathcal{P}\!\mathit{ic}}
\def\Picardfunctor{\mathrm{Pic}}
\def\Deformationcategory{\mathcal{D}\!\mathit{ef}}
\usepackage{cleveref}
\crefname{originalreference}{original reference}{original references}
\Crefname{originalreference}{Original reference}{Original references}
\makeatletter
\newlabel{example-CM-irreducible}{{0F91}{source}{}{sourcepointer-1}{}}
\newlabel{example-CM-irreducible@cref}{{[originalreference][0][]0F91}{[1][0][]source}{}{}{}}
\makeatother
\crefname{originallemma}{original lemma}{original lemmas}
\Crefname{originallemma}{Original lemma}{Original lemmas}
\makeatletter
\newlabel{local-cohomology-lemma-catenary-S2-equidimensional}{{0FIW}{source}{}{sourcepointer-2}{}}
\newlabel{local-cohomology-lemma-catenary-S2-equidimensional@cref}{{[originallemma][0][]0FIW}{[1][0][]source}{}{}{}}
\makeatother
\crefname{originallemma}{original lemma}{original lemmas}
\Crefname{originallemma}{Original lemma}{Original lemmas}
\makeatletter
\newlabel{properties-lemma-normal-locally-Noetherian}{{033N}{source}{}{sourcepointer-3}{}}
\newlabel{properties-lemma-normal-locally-Noetherian@cref}{{[originallemma][0][]033N}{[1][0][]source}{}{}{}}
\makeatother
\crefname{originalsection}{original section}{original sections}
\Crefname{originalsection}{Original section}{Original sections}
\makeatletter
\newlabel{section-setup}{{02QK}{source}{}{sourcepointer-4}{}}
\newlabel{section-setup@cref}{{[originalsection][0][]02QK}{[1][0][]source}{}{}{}}
\makeatother
\crefname{originalreference}{original reference}{original references}
\Crefname{originalreference}{Original reference}{Original references}
\makeatletter
\newlabel{situation-setup}{{02QL}{source}{}{sourcepointer-5}{}}
\newlabel{situation-setup@cref}{{[originalreference][0][]02QL}{[1][0][]source}{}{}{}}
\makeatother
\crefname{originallemma}{original lemma}{original lemmas}
\Crefname{originallemma}{Original lemma}{Original lemmas}
\makeatletter
\newlabel{topology-lemma-dimension-function-unique}{{02IB}{source}{}{sourcepointer-6}{}}
\newlabel{topology-lemma-dimension-function-unique@cref}{{[originallemma][0][]02IB}{[1][0][]source}{}{}{}}
\makeatother
\title{Equivalence of three local forms of local equidimensionality for the relative dimension function}
\author{The Stacks Project Authors}
\date{stacks-project:chow.tex:local-snapshot-20261012; GNU FDL 1.2 or later}
\begin{document}
\maketitle
\noindent\textit{Editorial note.} All delivered mathematical text is the author's own, copied from the pinned Stacks Project chapter file chow.tex; no statement or proof step was rewritten, reordered or removed.\par\medskip
\noindent\textit{Editorial note.} Preamble provenance: the pinned chapter file opens with \textbackslash{}input\{preamble\} and carries no preamble text of its own, so the project's single Stacks-wide preamble staged at sources/stacks-chow/preamble.tex is excerpted verbatim as three spans, lines 8-12 (the verbatim package together with the project's own reference, slogan and history display-environment definitions, including the author's commented-out older form of the reference environment), lines 14-18 (Xy-pic setup) and lines 154-217 (the native theorem declarations and mathematical macros). The lines 8-12 span is required because a counted statement or its surrounding text may use those display environments, which the project typesets as comments. Nothing was added, removed or redefined; the wrapper class and the amsmath/amssymb/amsthm/graphicx base packages are supplied by the assembler.\par\medskip
\noindent\textit{Editorial note.} Same-chapter pointers are located at their \textbackslash{}label lines inside the pinned chapter file; cross-chapter pointers are located in the project's own tag list staged in the snapshot (sources/stacks-chow/tags.txt), at the verbatim lines carrying their canonical tags.\par\medskip
\noindent\textit{Editorial note.} The author's footnote attached to the Cohen-Macaulay clause of the statement (source lines 8294-8296) is retained verbatim inside the counted statement.\par\medskip
\noindent\textit{Allowed standard prerequisites (editorial).}\begin{itemize}
\item A normal locally Noetherian scheme is a disjoint union of integral schemes. Properties, label properties-lemma-normal-locally-Noetherian, invoked at proof line 8307. Gives the normal case, since on each integral piece the dimension function is constant at generic points.
\item Uniqueness up to a locally constant function for dimension functions, and the Cohen-Macaulay example. Topology, label topology-lemma-dimension-function-unique (proof line 8313); Example example-CM-irreducible (proof line 8311). Shows that the negative of the dimension of the local ring is a dimension function which vanishes at generic points, so that the relative dimension function differs from it by a locally constant function, giving the Cohen-Macaulay case.
\end{itemize}\medskip
\ifdefined\hypertarget\hypertarget{sourcepointer-1}{}\fi
\noindent\textit{Source reference.} Original reference 0F91: sources/stacks-chow/pinned-original-source.tex, source lines 1433--1433. Navigation pointer for \textbackslash{}ref\{example-CM-irreducible\} used inside the counted span (line 8311); the argument that consumes it is retained in the span.\par\medskip
\ifdefined\hypertarget\hypertarget{sourcepointer-2}{}\fi
\noindent\textit{Source reference.} Original lemma 0FIW: sources/stacks-chow/tags.txt, source lines 19406--19406. Navigation pointer for \textbackslash{}ref\{local-cohomology-lemma-catenary-S2-equidimensional\} used inside the counted span (line 8296); the argument that consumes it is retained in the span.\par\medskip
\ifdefined\hypertarget\hypertarget{sourcepointer-3}{}\fi
\noindent\textit{Source reference.} Original lemma 033N: sources/stacks-chow/tags.txt, source lines 4062--4062. Navigation pointer for \textbackslash{}ref\{properties-lemma-normal-locally-Noetherian\} used inside the counted span (line 8307); the argument that consumes it is retained in the span.\par\medskip
\ifdefined\hypertarget\hypertarget{sourcepointer-4}{}\fi
\noindent\textit{Source reference.} Original section 02QK: sources/stacks-chow/pinned-original-source.tex, source lines 1357--1357. Navigation pointer for \textbackslash{}ref\{section-setup\} used inside the counted span (line 8280); the argument that consumes it is retained in the span.\par\medskip
\ifdefined\hypertarget\hypertarget{sourcepointer-5}{}\fi
\noindent\textit{Source reference.} Original reference 02QL: sources/stacks-chow/pinned-original-source.tex, source lines 1372--1372. Navigation pointer for \textbackslash{}ref\{situation-setup\} used inside the counted span (line 8278); the argument that consumes it is retained in the span.\par\medskip
\ifdefined\hypertarget\hypertarget{sourcepointer-6}{}\fi
\noindent\textit{Source reference.} Original lemma 02IB: sources/stacks-chow/tags.txt, source lines 3294--3294. Navigation pointer for \textbackslash{}ref\{topology-lemma-dimension-function-unique\} used inside the counted span (line 8313); the argument that consumes it is retained in the span.\par\medskip
\begin{lemma}
\label{lemma-locally-equidimensional}
Let $(S, \delta)$ be as in Situation \ref{situation-setup}.
Let $X$ be locally of finite type over $S$. Write
$\delta = \delta_{X/S}$ as in Section \ref{section-setup}.
The following are equivalent
\begin{enumerate}
\item There exists a decomposition $X = \coprod_{n \in \mathbf{Z}} X_n$
into open and closed subschemes such that $\delta(\xi) = n$ whenever
$\xi \in X_n$ is a generic point of an irreducible component of $X_n$.
\item For all $x \in X$ there exists an open neighbourhood $U \subset X$
of $x$ and an integer $n$ such that $\delta(\xi) = n$ whenever
$\xi \in U$ is a generic point of an irreducible component of $U$.
\item For all $x \in X$ there exists an integer $n_x$ such that
$\delta(\xi) = n_x$ for any generic point $\xi$ of an irreducible
component of $X$ containing $x$.
\end{enumerate}
The conditions are satisfied if $X$ is either
normal or Cohen-Macaulay\footnote{In fact, it suffices if
$X$ is $(S_2)$. Compare with Local Cohomology, Lemma
\ref{local-cohomology-lemma-catenary-S2-equidimensional}.}.
\end{lemma}

\begin{proof}
It is clear that (1) $\Rightarrow$ (2) $\Rightarrow$ (3).
Conversely, if (3) holds, then we set $X_n = \{x \in X \mid n_x = n\}$
and we get a decomposition as in (1). Namely, $X_n$ is open because
given $x$ the union of the irreducible components of $X$ passing through $x$
minus the union of the irreducible components of $X$ not passing through $x$
is an open neighbourhood of $x$. If $X$ is normal, then $X$ is a
disjoint union of integral schemes
(Properties, Lemma \ref{properties-lemma-normal-locally-Noetherian})
and hence the properties hold.
If $X$ is Cohen-Macaulay, then
$\delta' : X \to \mathbf{Z}$, $x \mapsto -\dim(\mathcal{O}_{X, x})$
is a dimension function on $X$ (see Example \ref{example-CM-irreducible}).
Since $\delta - \delta'$ is locally constant
(Topology, Lemma \ref{topology-lemma-dimension-function-unique})
and since $\delta'(\xi) = 0$ for every generic point $\xi$ of $X$
we see that (2) holds.
\end{proof}

\end{document}
